\documentclass[10pt,a4paper]{amsart}
\usepackage[margin=19mm]{geometry}
\usepackage{amsmath,amssymb,mathtools}
\usepackage[scr=rsfs,cal=euler]{mathalfa}
\usepackage{xfrac}
\usepackage[unicode,hidelinks,bookmarks=false]{hyperref}
\newcommand{\ie}{i.e.\ }
\newcommand{\cf}{cf.\ }
\newcommand{\resp}{respectively\ }
\newcommand{\Subsection}{\subsection}
\providecommand{\nobreakdash}{\nobreak\hbox{-}}
\setlength{\emergencystretch}{2em}
\multlinegap0pt
\usepackage[shortlabels]{enumitem}
\usepackage{mathtools}
\usepackage{amssymb}
\newtheorem{definition}{Definition}
\newtheorem{theorem}{Theorem}
\newcommand{\N}{\mathbb N}
\title{A scalar $c_0$-approximation criterion and gap-coordinate preduals for Lipschitz-free spaces}
\author{Luigi D'Onofrio}
\date{Source: arXiv:2606.22577v1 (CC BY 4.0)}
\begin{document}
\maketitle

\noindent\emph{Source and licence.} A scalar $c_0$-approximation criterion and gap-coordinate preduals for Lipschitz-free spaces. Version arXiv:2606.22577v1 (CC BY 4.0), distributed by arXiv under the Creative Commons Attribution 4.0 International licence, http://creativecommons.org/licenses/by/4.0/, as recorded in the arXiv licence metadata for this exact version and shown on the abstract page https://arxiv.org/abs/2606.22577v1. Full licence notice: \texttt{licenses/arxiv-2606.22577-CC-BY-4.0.txt}.\par\smallskip

\noindent\textit{Editorial scope.} Counted unit: the article's theorem thm:scalar-c0 (source0.tex lines 230--260). The article's complete original proof of it is delivered with it (source0.tex lines 262--384, one \texttt{proof} environment of 123 lines). No mathematics is written, repaired or re-derived: the statement and the proof are the article's own lines, in the article's own notation, and no retained line is substituted or re-flowed.  Retained uncounted as the source-local context the proof needs: lines 185--217.  Nothing was omitted from any retained span, so the package carries no omission record.  No retained line contains a citation, so the wrapper carries no bibliography and no bibliography entry.  No retained line contains a figure environment, an image-inclusion command or drawing code, so no image is delivered and the package declares no asset dependency.  Editorial formatting only, in addition: an AMS \texttt{amsart} preamble that restates the article's own declarations and macros used by the retained lines, namely the environments \texttt{definition}, \texttt{theorem} on separate counters, together with the standard packages those lines need.  The retained environments print in this document as Definition 1, Theorem 1 in order of appearance, whereas the article prints the same items with the numbering of its own full document.  The complete native source of the article as distributed in the arXiv e-print for this exact version is bundled unchanged as \texttt{sources/203359/source0.tex}.
\begin{definition}\label{def:scalar-setting}
Let $E$ be a Banach space and let $(f_n)_{n\in\N}\subset E^*$ satisfy
\[
        \sup_{n\in\N}\|f_n\|\le 1,
        \qquad
        \|x\|=\sup_{n\in\N}|f_n(x)|\quad (x\in E).
\]
Define
\[
        V:E\longrightarrow \ell^\infty,
        \qquad
        Vx=(f_n(x))_{n\in\N}.
\]
Then $V$ is an isometric embedding. Define also
\[
        V^\sharp:\ell^1\longrightarrow E^*,
        \qquad
        V^\sharp\xi=\sum_{n=1}^\infty \xi_n f_n.
\]
Equivalently, $V^\sharp$ is the restriction of the adjoint
$V^*:(\ell^\infty)^*\to E^*$ to the canonical subspace
$\ell^1\subset (\ell^\infty)^*$.

Let
\[
        Y:=V^\sharp(\ell^1)\subset E^*,
\]
e equip $Y$ with the quotient norm inherited from $\ell^1/\ker V^\sharp$.
Finally, set
\[
        E_{c_0}:=\{x\in E:Vx\in c_0\}.
\]
\end{definition}

\begin{theorem}[Scalar $c_0$-criterion]\label{thm:scalar-c0}
Assume, in the setting of Definition~\ref{def:scalar-setting}, that
\[
        J:E\longrightarrow Y^*
\]
is an isometric isomorphism. Then the following assertions are equivalent.

\begin{enumerate}[label=\textup{(\roman*)}]
\item The canonical restriction map
\[
        \Lambda:Y\longrightarrow E_{c_0}^*,
        \qquad
        \Lambda(y)=y|_{E_{c_0}},
\]
is an isometric isomorphism.

\item For every $x\in E$ there exists a sequence $(x_k)\subset E_{c_0}$
such that
\[
        \sup_k\|x_k\|<\infty
        \qquad\text{and}\qquad
        f_n(x_k)\longrightarrow f_n(x)
        \quad(n\in\N).
\]
\end{enumerate}

When these conditions hold, the sequence in \textup{(ii)} may be chosen
with $\sup_k\|x_k\|\le \|x\|$. Moreover, on bounded sequences in $E$,
coordinatewise convergence on the family $(f_n)$ is equivalent to
convergence in $\sigma(E,Y)$.
\end{theorem}

\begin{proof}
Set
\[
        W:=V(E_{c_0})=V(E)\cap c_0\subset c_0.
\]
The restriction $V:E_{c_0}\to W$ is an isometric isomorphism. Hence
\[
        E_{c_0}^*\cong W^*.
\]
Since $W$ is a closed subspace of $c_0$, the classical duality of subspaces
of $c_0$ gives
\[
        W^*\cong \ell^1/W^\perp,
\]
where
\[
        W^\perp=\left\{\xi\in\ell^1:
        \sum_{n=1}^\infty \xi_n w_n=0\text{ for every }w=(w_n)\in W
        \right\}.
\]
On the other hand, by definition of the quotient norm on $Y$,
\[
        Y\cong \ell^1/\ker V^\sharp.
\]
We first observe that
\[
        \ker V^\sharp\subset W^\perp.
\]
Indeed, if $\xi\in\ker V^\sharp$ and $w=Vx\in W$ with $x\in E_{c_0}$, then
\[
        \sum_{n=1}^\infty \xi_n w_n
        =\sum_{n=1}^\infty \xi_n f_n(x)
        =V^\sharp\xi(x)=0.
\]
Therefore the canonical quotient map
$\ell^1/\ker V^\sharp\to \ell^1/W^\perp$ is contractive; it is precisely
$\Lambda:Y\to E_{c_0}^*$ under the identifications above.

Assume now \textup{(ii)}. Let $\xi\in W^\perp$ and let $x\in E$. Choose
$(x_k)\subset E_{c_0}$ as in \textup{(ii)}, say $\sup_k\|x_k\|\le M$. For
every $k$,
\[
        0=\sum_{n=1}^\infty \xi_n f_n(x_k),
\]
because $Vx_k\in W$ and $\xi$ annihilates $W$. Since
\[
        |\xi_n f_n(x_k)|\le M|\xi_n|,
        \qquad \xi=(\xi_n)\in \ell^1,
\]
the dominated convergence theorem for series yields
\[
        \sum_{n=1}^\infty \xi_n f_n(x)
        =\lim_{k\to\infty}\sum_{n=1}^\infty \xi_n f_n(x_k)=0.
\]
Thus $V^\sharp\xi=0$, and therefore $W^\perp\subset \ker V^\sharp$.
Consequently
\[
        W^\perp=\ker V^\sharp.
\]
It follows that
\[
        E_{c_0}^*\cong W^*\cong \ell^1/W^\perp
        =\ell^1/\ker V^\sharp\cong Y
\]
canonically and isometrically. This proves \textup{(i)}.

Conversely, assume \textup{(i)}. Then
$\Lambda:Y\to E_{c_0}^*$ is an isometric isomorphism, and hence its adjoint
\[
        \Lambda^*:E_{c_0}^{**}\longrightarrow Y^*
\]
is an isometric isomorphism. Let $x\in E$ and set
\[
        \widetilde x=(\Lambda^*)^{-1}(Jx)\in E_{c_0}^{**}.
\]
Then $\|\widetilde x\|=\|x\|$. By Goldstine's theorem, there exists a net
$(u_\alpha)\subset E_{c_0}$ such that
\[
        \sup_\alpha\|u_\alpha\|\le \|x\|,
        \qquad
        \kappa_{E_{c_0}}(u_\alpha)\xrightarrow{w^*}\widetilde x
        \quad\text{in }E_{c_0}^{**}.
\]
For $u\in E_{c_0}$ and $y\in Y$ one has
\[
        \langle \Lambda^*\kappa_{E_{c_0}}(u),y\rangle
        =\langle \kappa_{E_{c_0}}(u),\Lambda y\rangle
        =(\Lambda y)(u)=y(u)=\langle Ju,y\rangle.
\]
Hence $\Lambda^*\kappa_{E_{c_0}}(u)=Ju$. Applying $\Lambda^*$ to the
convergent net gives
\[
        Ju_\alpha\xrightarrow{\sigma(Y^*,Y)}Jx.
\]

Since $Y$ is a quotient of $\ell^1$, it is separable. Let $(y_j)_{j\ge1}$
be dense in the unit ball of $Y$. We choose a sequence from the net as
follows. For each $k$ choose an index $\alpha_k$ such that
\[
        |(Ju_{\alpha_k}-Jx)(y_j)|<\frac1k,
        \qquad 1\le j\le k.
\]
This is possible because the net converges to $Jx$ in $\sigma(Y^*,Y)$.
Set $x_k=u_{\alpha_k}$. Then $\sup_k\|x_k\|\le\|x\|$. If $y\in Y$ and
$\varepsilon>0$, choose $j$ such that $\|y-y_j\|$ is sufficiently small;
using the uniform boundedness of $(Jx_k)$ and $Jx$, one obtains
\[
        (Jx_k)(y)\longrightarrow (Jx)(y).
\]
Thus $x_k\to x$ in $\sigma(E,Y)$. Since each $f_n$ belongs to $Y$, this
implies
\[
        f_n(x_k)\longrightarrow f_n(x)\qquad(n\in\N),
\]
and proves \textup{(ii)}.

It remains to justify the final assertion. The linear span of $(f_n)$ is
dense in $Y$, because finitely supported sequences are dense in $\ell^1$ and
$Y$ is the quotient image of $\ell^1$. If $(z_k)$ is bounded in $E$ and
$f_n(z_k)\to f_n(z)$ for every $n$, then convergence holds first for finite
linear combinations of the $f_n$ and then for every $y\in Y$ by density and
boundedness. Thus $z_k\to z$ in $\sigma(E,Y)$. The converse is immediate.
\end{proof}

\end{document}
